\documentclass[10pt,a4paper]{article}
\usepackage{fontspec,amsmath,amssymb,graphicx}
\usepackage[margin=18mm]{geometry}
\setmainfont{Times New Roman}
\setlength{\parindent}{0pt}
\setlength{\parskip}{5pt}
\allowdisplaybreaks
\sloppy
\emergencystretch=2em
\begin{document}
\section*{BT\_3\_4\_3\_04 — Lời giải bổ sung}
\subsection*{Kiến thức và trình tự giải}
\textbf{Bước 1 — Đưa điều kiện biên về thuần nhất.} Với $u(0,t)=\mu(t)$, $u(L,t)=\nu(t)$, đặt $b=(1-x/L)\mu+(x/L)\nu$. Với Neumann hai đầu đặt $b=\mu x+(\nu-\mu)x^2/(2L)$. Với Dirichlet–Neumann đặt $b=\mu+x\nu$. Viết $u=b+w$; khi đó $w$ có biên thuần nhất, nguồn $Q=F+a^2b_{xx}-b_{tt}$ và dữ liệu $f-b(x,0)$, $g-b_t(x,0)$.

\textbf{Bước 2 — Chiếu nguồn và dữ liệu lên cùng hệ hàm riêng.} Tính $Q_n(t)=\int QX_n/\int X_n^2$, $A_n$, $B_n$. Không bỏ các dữ liệu mới phát sinh từ hàm nâng biên; chúng có thể khác không ngay khi dữ liệu gốc bằng không.

\textbf{Bước 3 — Giải ODE có nguồn cho từng mode.}


\[
y_n''+\omega_n^2y_n=Q_n(t),\qquad
y_n=A_n\cos\omega_nt+\frac{B_n}{\omega_n}\sin\omega_nt+\frac1{\omega_n}\int_0^t\sin(\omega_n(t-s))Q_n(s)ds.
\]


Với mode hằng Neumann, thay công thức bằng $y_0=A_0+B_0t+\int_0^t(t-s)Q_0(s)ds$.

\textbf{Bước 4 — Xử lý cộng hưởng trước khi chia.} Nguồn $\cos\alpha t$ có đáp ứng dữ liệu không là $(\cos\alpha t-\cos\omega t)/(\omega^2-\alpha^2)$; tại $\alpha=\omega$ lấy $t\sin\omega t/(2\omega)$. Nguồn $\sin\alpha t$ có đáp ứng $(\sin\alpha t-(\alpha/\omega)\sin\omega t)/(\omega^2-\alpha^2)$; tại cộng hưởng là $\sin\omega t/(2\omega^2)-t\cos\omega t/(2\omega)$. Không thế vào phân số $0/0$.

\textbf{Bước 5 — Cộng lại hàm nâng và kiểm tra duy nhất.} Thế $u=b+\sum y_nX_n$ vào PDE và mọi dữ liệu. Với hiệu hai nghiệm, nguồn và biên bằng không; nhân phương trình với $w_t$, tích phân từng phần để có $E'=0$, $E(0)=0$, suy $w=0$. Các góc không tương thích phải được ghi riêng.

\textbf{Chi tiết xác nhận tích phân nguồn.} Với $D(t)=\omega^{-1}\int_0^t\sin(\omega(t-s))Q(s)ds$, lần đạo hàm đầu cho $D'=\int_0^t\cos(\omega(t-s))Q(s)ds$, vì $\sin0=0$. Lần sau cho $D''=Q(t)-\omega\int_0^t\sin(\omega(t-s))Q(s)ds=Q-\omega^2D$. Do $D(0)=D'(0)=0$, tích phân này thêm đúng nguồn mà giữ nguyên hai dữ liệu đầu.
\subsection*{Lời giải áp dụng cho bài}
\textbf{Bước 1. Nhận dạng một mode.} Biên cố định trên $(0,1)$ và nguồn $\sin\pi x\sin t$ gợi $u=\sin\pi x\,T(t)$. Đạo hàm $u_{xx}=-\pi^2\sin\pi x\,T$; vì tốc độ bằng $1$, PDE trở thành $T''+\pi^2T=\sin t$, dữ liệu $T(0)=T'(0)=0$.

\textbf{Bước 2. Nghiệm riêng và thuần nhất.} Nguồn tần số $1$ khác $\pi$, nên thử $T_p=C\sin t$, được $C=1/(\pi^2-1)$. Nghiệm tổng $T=C\sin t+A\cos\pi t+B\sin\pi t$. Dữ liệu đầu cho $A=0$, $C+\pi B=0$, nên $B=-1/(\pi(\pi^2-1))$.

\textbf{Bước 3. Nghiệm.}
\[u(x,t)=\frac{\sin t-\sin\pi t/\pi}{\pi^2-1}\sin\pi x.\]
\textbf{Bước 4. Kiểm tra.} Hai biên bằng không vì $\sin\pi x=0$; $T(0)=0$ và $T'(0)=(1-1)/(\pi^2-1)=0$. Thế $T''+\pi^2T$ cho $\sin t$, nên đúng nguồn. Tính duy nhất theo năng lượng của hiệu hai nghiệm biên cố định.\end{document}